\paragraph{The finite analytic input in more detail.}
Here is a reconstruction of the estimates behind the three inherited facts. For a rational edge $[p,q]$ set $z(t)=p+t(q-p)$ and use left Riemann sums on a dyadic subdivision. If the integrand varies by at most $2^{-k}$ on each panel, comparing a sum with any dyadic refinement costs at most $|q-p|_1 2^{-k}$, plus the length times the sum of the two point-evaluation errors for independently sampled sums. For a continuity modulus $\nu$, taking $2^m\ge |q-p|_1 2^{\nu(k)}$ ensures the required panel size. Sorting the union of two rational node lists gives a common refinement, also for a split at a nondyadic rational parameter, and the same oscillation bound applies. This proves additivity and orientation reversal as name equalities.

For a triangle with taxicab diameter $D$ and perimeter $P$, subdivide repeatedly by the three midpoints. At level $m$ there are $4^m$ triangles of diameter at most $D2^{-m}$ and perimeter $P2^{-m}$. Choose this diameter below the input Taylor radius at $k$. On each triangle subtract the affine function $f(x)+d(x)(z-x)$ at a fixed vertex $x$. Its integral on the triangle boundary is zero, by integrating the constant and linear terms along the three edges, or by telescoping their quadratic primitive. The remaining integrand has norm at most $D2^{-m}2^{-k}$. All internal edges cancel, and the total integral has norm at most $DP2^{-k}$. For finite sampled sums the total perimeter at this level is $2^mP$, so extra error $\varepsilon$ is achieved by per-unit-length error $\varepsilon/(1+2^mP)$. This proves the finite-chain Cauchy identity by an arbitrary rational tolerance, including signed multiplicities.

For the square kernel, define $\pi=4\int_0^1(1+t^2)^{-1}\,dt$; the integral has a PR modulus from a rational Lipschitz bound and $\pi\ge2$. On the right edge of the unit square the integrand $dz/z$ is $(t+i)\,dt/(1+t^2)$ for $-1\le t\le1$. The odd part cancels and rotation gives the same value on all four sides. Hence the total is $2\pi i$, at any positive scale by change of variables. For $z^j$ with integer $j\ne-1$, $z^{j+1}/(j+1)$ is a primitive on the contour; the usual telescoping edge proof follows from the finite binomial expansion for positive powers, and from the difference of inverses under the edge's positive separation for negative ones.

For completeness, the annular Cauchy formula at a rational $v\in\partial Q_r(0)$ uses the kernel $g(u)/(u-v)$. Exclude a rational square of half-side $h$ about $v$ lying between $Q_s$ and $Q_R$. The remaining square region with two disjoint square holes is triangulated by drawing horizontal and vertical lines through all sides, retaining its rectangles and taking two triangles in each. Division by $u-v$ there is separated: multiplication and the inverse difference identity supply the required moduli. The Cauchy identity equates outer minus inner integrals to the integral on the small square. Replacing $g(u)$ by $g(v)$ there has integral error at most $(8h)(2/h)2^{-k}$ once $2h\le2^{-\nu(k)}$. The square moment then gives $2\pi i g(v)$ to arbitrary precision. Finally expand the two kernels by the \emph{finite} geometric identity with remainder (on the inner square expand in $u/v$, on the outer square in $v/u$). The explicit bound below results from their $1/2$ ratios. Thus the inherited analytic argument uses rational quadrature and the stated modulus of differentiability, even when $v$ is on an intermediate square edge. These proofs still require deep formalization for their use in ordinary iPRA.
